\section{La lógica del aprendizaje robótico autónomo}
\subsection{La brecha entre Human-in-the-loop y Autonomy}
El sistema actual de aprendizaje robótico depende de la supervisión humana: la función de recompensa, el reinicio del entorno y la recolección de datos están a cargo de ingenieros. La clase «Autonomy» plantea un reto: que el robot explore activamente el complejo mundo físico como lo hace un bebé, en lugar de esperar a que una persona le fije los objetivos.

\begin{figure}[H]
\centering
\includegraphics[width=0.92\textwidth]{slides-images/slide-01.jpg}
\caption{La Slide 1 presenta el tema central de la clase: Autonomy es, a la vez, el conjunto de la configuración de reward, reset y goal.}
\end{figure}

\begin{importantbox}{Los tres objetivos de Autonomy}
1. Construir la recompensa de forma automática o entender el resultado de la tarea;\\
2. Operar de forma continua en un mundo sin reinicios;\\
3. Elegir objetivos de forma autónoma y explorar con seguridad.
\end{importantbox}

\subsection{Criterios de evaluación frente a la realidad}
El ponente enfatiza tres dimensiones de evaluación: \textbf{capacidad de aprendizaje} (si se puede obtener suficiente signal), \textbf{escalabilidad} (si no depende de reinicios manuales) y \textbf{sostenibilidad} (si puede funcionar a largo plazo y transferirse a objetivos nuevos). La lógica de la exposición: se parte de la reward y, después, a partir de la lógica del entorno y de los objetivos, se avanza hacia la seguridad y la integración de sistemas.

\begin{knowledgebox}{Mapa de la lógica de la exposición}
La línea de investigación de Autonomy sigue este orden: \textbf{recompensa automática}→\textbf{control Reset-free}→\textbf{Goal Management}→\textbf{Safe Exploration}→\textbf{monitoreo del sistema}. Cada etapa parte de un supuesto anterior (reward, reset, goal) y describe cómo retirarlo.
\end{knowledgebox}

\subsection{Resumen de la sección}
Esta sección repasa la situación actual y la lógica de la exposición del aprendizaje robótico autónomo, y toma las cinco etapas «recompensa-entorno-objetivo-seguridad-sistema» como marco rector de las secciones siguientes.
\newpage
